%! Author = chtompki
%! Date = 1/14/26
\documentclass[12pt]{amsart}
% Default margins are too wide all the way around. I reset them here
\setlength{\hoffset}{-.75in}
\setlength{\textwidth}{6.5in}
\setlength{\voffset}{-.5in}
\setlength{\textheight}{9.0in}
\usepackage{amsmath}
\usepackage{leftindex}
\usepackage{amssymb}
\usepackage{amsthm}
\usepackage{enumitem}
\usepackage{setspace}

\newenvironment{solution}
{\begin{proof}[Solution]}
{\end{proof}}

\title{Combinatorics\\Final Exam}
\author{Rob Tompkins}
\date{\today}

\begin{document}
    \maketitle
    \date{\today}
    \begin{onehalfspace}
        % Problem 1
        \noindent\textbf{1.} (Multinomial Coefficients) A committee of 12 people is to be divided into
        four subcommittees of sizes 3, 4, 2, and 3.
        \begin{itemize}
            \item[(a)] How many ways can this be done if the four subcommittees have distinct responsibilities?
            \item[(b)] How many ways can this be done if two of the subcommittees of size 3 have identical
            responsibilities, while the others are distinct?
            \item[(c)] Explain why the multinomial coefficient
            \[
                \binom{12}{3.4.2.3}
            \]
            appears naturally in part (a).
        \end{itemize}

        \begin{solution}
        \end{solution}

        % Problem 2
        \noindent\textbf{2.} (Binomial Identies) Prove the identity
        \[
            \sum_{k=0}^n \binom{n}{k}^2 = \binom{2n}{k}.
        \]
        Give \textbf{two different proofs}
        \begin{itemize}
            \item[(a)] An algebraic proof using the binomial theorem.
            \item[(b)] A combinatorial proof by counting the same set in two ways.
        \end{itemize}

        \begin{proof}
        \end{proof}

        % Problem 3
        \noindent\textbf{3.} (Ramnsey Numbers)
        \begin{itemize}
            \item[(a)] Show that if $m\ge 3$ and $n\ge 3$ are integers then
            \[
                R(m,n) \le R(m-1,n) + R(m.n-1).
            \]
            \item[(b)] Show that if $R(m-1,n)$ and $R(m,n-1)$ are both even and greater than 2,
            then
            \[
                R(m,n) \le R(m-1,n) + R(m,n-1) - 1.
            \]
            \item[(c)] Show that $R(4,3) = 9$, $R(53) = 14$, and $R(4,4)=18$
        \end{itemize}

        \begin{proof}
        \end{proof}

        % Problem 4
        \noindent\textbf{4.} (Stirling Numbers of the Second Kind) Let $S(n,k)$ denote the Stirling
        numbers of the second kind, the number of ways to partition an $n$-element set into $k$
        nonempty unlabeled blocks.
        \begin{itemize}
            \item[(a)] Prove the recurrence
            \[
                S(n,k) = kS(n-1,k) + S(n-1,k-1).
            \]
            \item[(b)] Use the recurrence to compute $S(5,2)$, $S(5,3)$, and $S(6,3)$.
            \item[(c)] Give a direct combinatorial interpretation of $S(6,3)$.
        \end{itemize}

        \begin{solution}
        \end{solution}

        % Problem 5
        \noindent\textbf{5.}

        \begin{solution}
        \end{solution}

        % Problem 6
        \noindent\textbf{6.}

        \begin{solution}
        \end{solution}

        % Problem 7
        \noindent\textbf{7.}

        \begin{solution}
        \end{solution}

        % Problem 8
        \noindent\textbf{8.}

        \begin{solution}
        \end{solution}

        % Problem 9
        \noindent\textbf{9.}

        \begin{solution}
        \end{solution}

        % Problem 10
        \noindent\textbf{10.}

        \begin{solution}
        \end{solution}

        % Problem 11
        \noindent\textbf{11.}

        \begin{solution}
        \end{solution}

        % Problem 12
        \noindent\textbf{12.}

        \begin{solution}
        \end{solution}

        % Problem 13
        \noindent\textbf{13.}

        \begin{solution}
        \end{solution}

        % Problem 14
        \noindent\textbf{14.}

        \begin{solution}
        \end{solution}

        % Problem 15
        \noindent\textbf{15.}

        \begin{solution}
        \end{solution}

    \end{onehalfspace}
\end{document}
